\chapter{Coding}

This chapter describes the key technologies used to implement Financial Fraud Detection Using Graph Machine Learning and presents representative source code from the most critical modules, including the rule engine, the graph analysis module, and the model training pipeline. The backend is entirely written in \textbf{Python}, with the analyst client in \textbf{JavaScript (React)}. The codebase is structured as a Python package under \code{backend/app/} with \code{requirements.txt} as the dependency manifest. This chapter introduces the primary languages and frameworks before presenting representative code snippets from the most critical modules. Every listing is taken from the application repository at commit \code{3800cc1}; where lines are left out, the omission is marked in the listing with a comment beginning \code{\# ...}.

\section{Python}
Python 3.10 was chosen for its:

\begin{itemize}
  \item Mature data and scientific ecosystem (pandas, NumPy, scikit-learn, XGBoost) which covers the entire feature engineering and modelling workload without leaving the language.
  \item First-class graph tooling through NetworkX, whose algorithm catalogue --- PageRank, centrality measures, community detection and cycle enumeration --- maps directly onto the structural patterns the project needs to detect.
  \item Strong typing via Pydantic v2 for request validation and response serialisation.
  \item Built-in \code{sqlite3}, \code{logging}, \code{datetime} and \code{bisect} standard library modules used extensively in the project.
\end{itemize}

Python's uniformity across the stack is a practical advantage here: the same language loads the dataset, runs the rules, builds the graph, trains the models and serves the API, so there is no serialisation boundary between the detection logic and the web layer.

\section{FastAPI}
FastAPI 0.111 provides the HTTP layer. Its key advantages in this project are:

\begin{itemize}
  \item \textbf{Pydantic integration}: Request bodies and query parameters are automatically validated; invalid payloads return a structured 422 response without any hand-written checks.
  \item \textbf{Dependency injection}: Both the database session and JWT authentication are implemented as FastAPI dependencies and applied declaratively to each protected route, so a router simply declares \code{Depends(get\_db)} in its signature.
  \item \textbf{Automatic OpenAPI documentation}: The interactive \code{/docs} interface is generated from the route signatures, which made the detection endpoints testable before the React client existed.
  \item \textbf{Router modularity}: Each resource lives in its own \code{APIRouter} module and is mounted under a common \code{/api} prefix, keeping the application factory small.
\end{itemize}

\section{React}
React 18, built with Vite 5, builds the analyst client. The application uses:

\begin{itemize}
  \item \textbf{Material UI 5} for the component library and theming, giving the data-dense tables, dialogs and chips a consistent appearance without custom CSS.
  \item \textbf{Axios} with a shared instance that holds the base URL and an interceptor which attaches the bearer token and redirects to the login page on a 401 response.
  \item \textbf{React Router 6} for client-side routing, with a protected-route wrapper guarding every analytical page.
  \item \textbf{Recharts} for the dashboard visualisations --- risk distribution, alert severity breakdown and transaction volume over time.
\end{itemize}

\section{Sample Coding}

\Needspace{30\baselineskip}
\subsection{rules.py --- Rule Engine}
\begin{lstlisting}[language=Python, caption={\code{rules.py} --- thresholds and the fan-in / fan-out rule (excerpt; omitted lines marked)}, label=lst:rules]
LARGE_AMOUNT = 100_000
RAPID_COUNT = 5
RAPID_WINDOW_MIN = 10
NEW_ACCOUNT_DAYS = 30
NEW_ACCOUNT_AMOUNT = 50_000
FAN_THRESHOLD = 8
FAN_WINDOW_H = 1
CIRCLE_WINDOW_H = 48

# ... RULE_LABELS, RULE_CODES and _load() omitted ...


def run_rules(txns: pd.DataFrame, accounts: pd.DataFrame) -> dict:
    """Return {rule_name: set(txn_uid)} for each rule."""
    hits = {name: set() for name in RULE_LABELS}
    if txns.empty:
        return hits
    txns = txns.sort_values("timestamp").reset_index(drop=True)

    # R1 large amount
    hits["large_amount"] = set(txns.loc[txns["amount"] > LARGE_AMOUNT, "txn_uid"])

    # ... R2 and R3 omitted ...

    # R4 fan-in / R5 fan-out (distinct counterparties within a window)
    fwin = timedelta(hours=FAN_WINDOW_H)

    def fan(group_col, counter_col, out_key):
        # Single pass per account: a Counter of the counterparties inside the
        # window gives the distinct count in O(1) per step, instead of slicing
        # the window and calling nunique() for every row.
        for _, g in txns.groupby(group_col):
            ts = g["timestamp"].tolist()
            others = g[counter_col].tolist()
            uids = g["txn_uid"].tolist()
            in_window = Counter()
            left = 0
            marked = 0  # uids[left:marked] are already in hits
            for right in range(len(ts)):
                in_window[others[right]] += 1
                while ts[right] - ts[left] > fwin:
                    in_window[others[left]] -= 1
                    if not in_window[others[left]]:
                        del in_window[others[left]]
                    left += 1
                if len(in_window) >= FAN_THRESHOLD:
                    hits[out_key].update(uids[max(left, marked):right + 1])
                    marked = right + 1

    fan("receiver_uid", "sender_uid", "fan_in")
    fan("sender_uid", "receiver_uid", "fan_out")

    # ... R6 omitted ...
    return hits
\end{lstlisting}

\subsection{graph\_analysis.py --- Graph Construction and Metrics}
\begin{lstlisting}[language=Python, caption={\code{graph\_build.py}, \code{graph\_analysis.py} and \code{fraud\_ring.py} --- construction, metrics and ring assembly (abridged excerpts)}, label=lst:graph]
# --- graph_build.py -------------------------------------------------------
def build_graph(db: Session) -> nx.DiGraph:
    G = nx.DiGraph()
    for uid, name, risk in db.query(Account.account_uid, Account.customer_name, Account.risk_level):
        G.add_node(uid, name=name, risk=risk)

    rows = db.query(
        Transaction.sender_uid, Transaction.receiver_uid, Transaction.amount,
    ).all()
    for sender, receiver, amount in rows:
        if G.has_edge(sender, receiver):
            e = G[sender][receiver]
            e["amount"] += amount
            e["count"] += 1
        else:
            G.add_edge(sender, receiver, amount=float(amount), count=1)
    return G


# --- graph_analysis.py ----------------------------------------------------
CYCLE_LENGTH_BOUND = 3


def analyze(db: Session) -> dict:
    G = build_graph(db)
    n, m = G.number_of_nodes(), G.number_of_edges()

    degree_centrality = nx.degree_centrality(G)
    pagerank = nx.pagerank(G, weight="amount")

    # ... weakly connected components omitted ...

    # Louvain communities (on the undirected projection)
    community_id, community_size = {}, {}
    communities = nx.community.louvain_communities(G.to_undirected(), seed=42)
    for i, com in enumerate(communities):
        for node in com:
            community_id[node] = i
            community_size[node] = len(com)

    # short directed cycles (money loops): 2-cycles and 3-cycles via set intersection.
    # (nx.simple_cycles is far too slow on ~80k edges; this is equivalent for length <= 3)
    succ = {u: set(G.successors(u)) for u in G.nodes}
    pred = {u: set(G.predecessors(u)) for u in G.nodes}
    cycle_nodes = set()
    for u, v in G.edges:
        if u < v and u in succ[v]:          # 2-cycle u <-> v (count once)
            cycle_nodes.update((u, v))
    for u, v in G.edges:
        common = succ[v] & pred[u]
        common.discard(u); common.discard(v)
        if common:
            cycle_nodes.add(u); cycle_nodes.add(v)
            cycle_nodes.update(common)

    # ... the cycle counters, money totals and the bulk insert into
    #     graph_metrics omitted ...


# --- fraud_ring.py --------------------------------------------------------
def detect(db: Session) -> dict:
    # ... per-account scores, names and labels, and the rules fired on each
    #     transaction (rules_by_txn), loaded from the database, omitted ...

    # 1) suspicious-activity subgraph; each edge remembers which rules fired on it
    G = nx.DiGraph()
    sus_txns = db.query(
        Transaction.txn_uid, Transaction.sender_uid, Transaction.receiver_uid, Transaction.amount,
    ).filter(Transaction.status == "suspicious").all()
    for txn_uid, s, r, amt in sus_txns:
        _add_edge(G, s, r, amt)
        G[s][r]["rules"] |= rules_by_txn.get(txn_uid, set())

    # ... transactions between accounts whose ML score is at least
    #     ML_SCORE_FLOOR (0.7) are added to G the same way, omitted ...

    # 2) connected groups; split oversized components into communities so
    #    distinct rings inside one big suspicious blob separate out
    MAX_COMPONENT = 60
    groups = []
    for comp in nx.weakly_connected_components(G):
        if len(comp) < MIN_RING_SIZE:
            continue
        if len(comp) <= MAX_COMPONENT:
            groups.append(comp)
            continue
        sub_u = G.subgraph(comp).to_undirected()
        for com in nx.community.louvain_communities(sub_u, seed=42, resolution=2.0):
            if len(com) >= MIN_RING_SIZE:
                groups.append(set(com))

    rings = []
    for comp in groups:
        sub = G.subgraph(comp)
        received = {u: 0.0 for u in comp}
        sent = {u: 0.0 for u in comp}
        for u, v, d in sub.edges(data=True):
            sent[u] += d["amount"]
            received[v] += d["amount"]

        # 3) roles
        main = max(comp, key=lambda u: received[u])
        withdrawal = None
        outs = sorted(G.successors(main), key=lambda v: G[main][v]["amount"], reverse=True)
        for v in outs:
            if v in comp and sent.get(v, 0) <= received.get(v, 0):
                withdrawal = v
                break
        # ... every other member is given the role "source" ...

        member_scores = [scores.get(u, 0.0) or 0.0 for u in comp]
        avg_score = sum(member_scores) / len(member_scores)
        ring_rules = set().union(*(d["rules"] for _, _, d in sub.edges(data=True)))
        rule_score = min(len(ring_rules), RULES_FOR_FULL_SCORE) / RULES_FOR_FULL_SCORE
        risk_score = max(avg_score, rule_score)
        # ... the ring record (members, roles, flow, rules, risk level) and
        #     the write to rings.json omitted ...
\end{lstlisting}

\noindent In \code{fraud\_ring.py}, \code{MIN\_RING\_SIZE} is 4 and \code{RULES\_FOR\_FULL\_SCORE} is 3, so one distinct rule on a ring's own transactions gives a rule score of 0.33, two give 0.67 and three or more give 1.0.

\subsection{baseline.py --- Model Training and Scoring}
The excerpt below is taken from \code{app/ml/baseline.py}; imports, the feature-importance summary and the write-back of scores to the database are omitted for space. Note that the three models handle class imbalance by three different mechanisms, and that Isolation Forest is unsupervised --- it is fitted on the training features alone and never sees a label. Its \code{contamination} parameter, however, is set from the label array of the whole dataset, test partition included; Section~\ref{sec:splits} discusses that choice and measures its effect.

\begin{lstlisting}[language=Python, caption={\code{baseline.py} --- model training and scoring (abridged excerpt)}, label=lst:baseline]
FEATURE_COLS = [
    "total_sent", "total_received", "txn_count_sent", "txn_count_received",
    "unique_receivers", "unique_senders", "avg_sent_amount", "avg_received_amount",
    "max_sent_amount", "max_received_amount", "txn_per_active_day",
    "in_degree", "out_degree", "degree_centrality", "pagerank", "in_cycle", "community_size",
]


def train_and_score(db: Session, test_size: float = 0.25, seed: int = 42) -> dict:
    df = load_features()
    if df is None:
        raise ValueError("Features not built. Run feature engineering first.")

    X = df[FEATURE_COLS].values
    y = df["is_fraud"].values
    X_train, X_test, y_train, y_test = train_test_split(
        X, y, test_size=test_size, stratify=y, random_state=seed,
    )
    n_pos, n_neg = int(y_train.sum()), int(len(y_train) - y_train.sum())

    results, models = {}, {}

    # --- Random Forest (class-weighted) ---
    rf = RandomForestClassifier(
        n_estimators=300, class_weight="balanced", random_state=seed, n_jobs=-1,
    )
    rf.fit(X_train, y_train)
    rf_score = rf.predict_proba(X_test)[:, 1]
    results["random_forest"] = _metrics(y_test, (rf_score >= 0.5).astype(int), rf_score)
    models["random_forest"] = rf

    # --- XGBoost (scale_pos_weight) ---
    xgb = XGBClassifier(
        n_estimators=300, max_depth=6, learning_rate=0.1,
        scale_pos_weight=n_neg / max(n_pos, 1),
        eval_metric="logloss", random_state=seed, n_jobs=-1,
    )
    xgb.fit(X_train, y_train)
    xgb_score = xgb.predict_proba(X_test)[:, 1]
    results["xgboost"] = _metrics(y_test, (xgb_score >= 0.5).astype(int), xgb_score)
    models["xgboost"] = xgb

    # --- Isolation Forest (unsupervised; contamination = observed fraud rate) ---
    iso = IsolationForest(
        n_estimators=200, contamination=float(y.mean()), random_state=seed, n_jobs=-1,
    )
    iso.fit(X_train)
    iso_pred = (iso.predict(X_test) == -1).astype(int)
    iso_score = -iso.score_samples(X_test)  # higher = more anomalous
    results["isolation_forest"] = _metrics(y_test, iso_pred, iso_score)

    # --- feature importance (from the best supervised model by F1) ---
    best_name = max(("random_forest", "xgboost"), key=lambda k: results[k]["f1"])
    best = models[best_name]
    # ... top-eight feature importances omitted ...

    # --- score ALL accounts with the best model, write back ---
    all_scores = best.predict_proba(X)[:, 1]
    # ... write-back to Account.fraud_score, the saved model and
    #     baseline_results.json omitted ...
\end{lstlisting}

\subsection{gnn.py --- Graph Neural Network Models}
Both models share one two-layer architecture, differing only in the convolution class. The layers are used with PyTorch Geometric's defaults and are trained full-batch over the whole graph: no neighbour sampling is configured, and \code{SAGEConv} therefore aggregates a node's neighbours by their mean and adds a separately learned transformation of the node's own features. Neither layer learns a weight for an individual neighbour.

\begin{lstlisting}[language=Python, caption={\code{gnn.py} --- model definition and weighted training (excerpt; omitted lines marked)}, label=lst:gnn]
def _build_model(model_type: str, in_dim: int, hidden: int = 64, out_dim: int = 2):
    import torch
    from torch_geometric.nn import GCNConv, SAGEConv

    conv_cls = {"gcn": GCNConv, "graphsage": SAGEConv}[model_type]

    class Net(torch.nn.Module):
        def __init__(self):
            super().__init__()
            self.conv1 = conv_cls(in_dim, hidden)
            self.conv2 = conv_cls(hidden, out_dim)
            self.dropout = torch.nn.Dropout(0.3)

        def forward(self, x, edge_index):
            x = self.conv1(x, edge_index).relu()
            x = self.dropout(x)
            return self.conv2(x, edge_index)

    return Net()


def train(db: Session, source: str = "synthetic", model_type: str = "graphsage",
          epochs: int = 100, hidden: int = 64, lr: float = 0.01, seed: int = 42) -> dict:
    # ... data assembly and the train / validation / test masks omitted ...

    # class-imbalance handling: weighted cross-entropy
    y_train = data.y[train_mask]
    counts = torch.bincount(y_train, minlength=2).float()
    weights = counts.sum() / (2 * counts.clamp(min=1))

    model = _build_model(model_type, data.num_features, hidden=hidden)
    optimizer = torch.optim.Adam(model.parameters(), lr=lr, weight_decay=5e-4)
    loss_fn = torch.nn.CrossEntropyLoss(weight=weights)

    # ... training-history bookkeeping omitted ...
    y_all = data.y
    for epoch in range(epochs):
        model.train()
        optimizer.zero_grad()
        out = model(data.x, data.edge_index)
        loss = loss_fn(out[train_mask], y_all[train_mask])
        loss.backward()
        optimizer.step()
        # ... periodic history recording omitted ...
\end{lstlisting}

The weights are the inverse class frequencies of the training partition. They are modest here rather than extreme: at the account level roughly a third of the seeded accounts carry a positive label, so the weighting corrects a moderate imbalance rather than a rare-event one. The severe imbalance in this dataset is at the transaction level, where only 3.33\% of transactions are fraudulent --- but the graph neural networks classify accounts, not transactions, so it is the account-level rate that applies to this loss.
